% SHARD-ID: AQM-25-QUASILOCAL-ZARISKI-ALGEBRA
% SHARD-TITLE: Quasi-Local Zariski Algebra
% SHARD-SUMMARY: Records the directed-system choice for building a quasi-local algebra from Zariski regions.
% SHARD-SUMMARY: Proves that open-local algebras over quasi-compact opens form the natural filtered system.
% SHARD-SUMMARY: Reinterprets the closed-complement construction as a closed-support system rather than the primary quasi-local net.
% SHARD-KEYWORDS: quasi-local algebra, directed system, filtered colimit, Zariski open, quasi-compact open, closed support

\section{Quasi-Local Zariski Algebra}
\label{sec:quasilocal-zariski-algebra}

\subsection{Status and Source Anchors}

This shard records the answer to the directed-system question left by
AQM-24.  The Zariski-topology source remains
\path{references/algebraic_geometry/stacks_project_algebra.tex}: standard
opens form a basis in lines 3104--3117, \(\operatorname{Spec}R\) is
quasi-compact in lines 3251--3273, and the topology has a quasi-compact open
basis with quasi-compact intersections in lines 3275--3298.

No analytic completion convention is fixed here.  The object constructed below
is the algebraic quasi-local \(*\)-algebra, i.e. the filtered colimit of the
local finite-region \(*\)-algebras.  A later \(C^*\)-completion should be
recorded as an additional convention.

\subsection{Lamport Dependency Skeleton}

\[
\begin{array}{rcl}
D1&:&\text{A directed set is a preorder in which any two elements have
      a common upper bound.}\\
D2&:&\mathsf O(X)=\{\text{Zariski opens of }X\},\quad
      \mathsf K(X)=\{\text{quasi-compact Zariski opens of }X\}.\\
D3&:&U\in\mathsf K(X)\mapsto\mathcal A(U)
      \text{ is the open-local algebra of AQM-24.}\\
L1&:&\mathsf O(X)\text{ is directed under inclusion.}\\
L2&:&\mathsf K(X)\text{ is directed under inclusion.}\\
T1&:&(\mathcal A(U))_{U\in\mathsf K(X)}
      \text{ is the natural quasi-local directed system.}\\
T2&:&\text{Closed complements are naturally indexed by closed supports.}
\end{array}
\]

\subsection{Directed Region Posets D1--L2}

A preorder \(I\) is called directed if for every \(i,j\in I\) there exists
\(k\in I\) such that \(i\leq k\) and \(j\leq k\).  A directed system of
\(*\)-algebras over \(I\) is a family \((A_i)_{i\in I}\) with structure maps
\(\alpha_{ij}:A_i\to A_j\) for \(i\leq j\), satisfying
\[
  \alpha_{ii}=\operatorname{id},\qquad
  \alpha_{jk}\alpha_{ij}=\alpha_{ik}.
\]

Let \(X=\operatorname{Spec}R\).  Write \(\mathsf O(X)\) for the set of
Zariski open subsets ordered by inclusion.  Write \(\mathsf K(X)\) for the
quasi-compact open subsets ordered by inclusion.  At the present concrete
level, \(\mathsf K(X)\) can be represented by finite unions of standard opens
\[
  D(f_1)\cup\cdots\cup D(f_n),
\]
using the Stacks source cited above.

\paragraph{Lemma \(L1\).}
\(\mathsf O(X)\) is directed under inclusion.

\emph{Proof.}
If \(U,V\subset X\) are open, then \(U\cup V\) is open and contains both.
Thus \(U\cup V\) is a common upper bound.

\paragraph{Lemma \(L2\).}
\(\mathsf K(X)\) is directed under inclusion.

\emph{Proof.}
If \(U,V\) are quasi-compact opens, then \(U\cup V\) is open.  To show it is
quasi-compact, let \((W_a)\) be an open cover of \(U\cup V\).  It restricts to
an open cover of \(U\) and to an open cover of \(V\).  Since \(U\) and \(V\)
are quasi-compact, finitely many \(W_a\)'s cover \(U\), and finitely many
\(W_a\)'s cover \(V\).  The union of those two finite subfamilies covers
\(U\cup V\).  Hence \(U\cup V\in\mathsf K(X)\), and it is a common upper
bound.

\subsection{The Quasi-Local System T1}

For \(U\in\mathsf K(X)\), let \(\mathcal A(U)\) be the open-local Weyl
observable algebra from AQM-24.  For \(U\subset V\), use the inclusion
\[
  \iota_{U,V}:\mathcal A(U)\hookrightarrow\mathcal A(V),\qquad
  a\mapsto a\otimes\mathbf 1_{V\setminus U}.
\]

\paragraph{Theorem \(T1\).}
\((\mathcal A(U),\iota_{U,V})_{U\in\mathsf K(X)}\) is a directed system of
\(*\)-algebras.  Its algebraic filtered colimit
\[
  \mathcal A_{\rm qloc}(X)
  :=
  \varinjlim_{U\in\mathsf K(X)}\mathcal A(U)
\]
is the natural Zariski quasi-local algebra for the present programme.

\emph{Proof.}
The index category is directed by Lemma \(L2\).  For every inclusion
\(U\subset V\), AQM-24 proves that \(\iota_{U,V}\) is an injective unital
\(*\)-homomorphism.  The same lemma proves functoriality:
\[
  \iota_{V,W}\circ\iota_{U,V}=\iota_{U,W}
\]
for \(U\subset V\subset W\).  These are exactly the axioms of a directed
system.  The displayed \(\varinjlim\) is therefore well-defined as the
algebraic colimit of this system.

\subsection{Closed Supports T2}

The closed-complement assignment from AQM-24 sends an open \(U\) to
\[
  \mathcal A(X\setminus U).
\]
With opens ordered by inclusion, this assignment is contravariant.  Therefore
it is not the primary quasi-local system if ``local region grows'' is meant to
match \(U\subset V\).

The same construction becomes covariant if it is indexed directly by closed
supports.  Let \(\mathsf Z(X)\) be the set of Zariski closed subsets ordered by
inclusion.  Since a finite union of closed subsets is closed, \(\mathsf Z(X)\)
is directed: \(Z,Z'\subset Z\cup Z'\).  The assignment
\[
  Z\longmapsto\mathcal A(Z)
\]
with the same tensor-by-identity inclusions is then a directed system.  This is
best read as a closed-support or vanishing-locus quasi-local algebra, not as
the open-local observable algebra.

\subsection{Finite Checks}

For \(X=\operatorname{Spec}\mathbb F_p\), the space has a largest open
\(X\).  Hence the directed system has a terminal object and
\[
  \mathcal A_{\rm qloc}(X)=\mathcal A(X)\cong M_p(\mathbb C).
\]

For \(X=\operatorname{Spec}(\mathbb Z/6\mathbb Z)\), all opens are
quasi-compact and \(X\) is again terminal.  The quasi-local algebra is
\[
  \mathcal A_{\rm qloc}(X)
  =
  M_2(\mathbb C)\otimes M_3(\mathbb C)\cong M_6(\mathbb C).
\]
The proper open-local subalgebras are the two commuting singleton algebras
\[
  M_2(\mathbb C)\otimes\mathbf 1_3,\qquad
  \mathbf 1_2\otimes M_3(\mathbb C),
\]
and their common upper bound is the full open \(X\).

\subsection{Conclusion}

For the quasi-local algebra, the natural general convention is
\[
  \boxed{\mathcal A_{\rm qloc}(X)
  =\varinjlim_{U\in\mathsf K(X)}\mathcal A(U)}
\]
with \(\mathsf K(X)\) the quasi-compact Zariski opens and with inclusions
following open inclusion.  The closed-complement possibility is retained as a
separate closed-support construction, useful for vanishing loci and defects,
but not as the primary open-local quasi-local algebra.
